\begin{abstract}
Humans learn an unfamiliar physical world by acting in it, and then see new scenes in the light of what they learned. We ask whether interaction gives a multimodal agent the same ability: after budgeted experiments in a physical world whose laws and constants may depart from canonical physics, does the agent understand that world's principles well enough to foresee how a \emph{new} scene from it will look? Improving on the practiced task is not enough, since an agent can do so by searching locally, replaying a trajectory that worked, or exploiting visual regularities of the scene it acted on. We introduce a benchmark built around this question. Each latent world is indexed by a physical-phenomenon family, a world-parameter regime, a latent-dynamics condition, and a latent-law realization, and each realization contains several counterfactual prediction instances. During support, the agent experiments on some of these instances through a tool interface that returns observable state but withholds hidden constants, the governing law, and answer keys. Tool access is then frozen. On held-out instances of the same realization, the agent must select the physically correct future frame from four rendered candidates whose positions are independently permuted. The benchmark contains 5,670 instances across seven phenomenon families. A support-budget sweep, a cross-realization control that pairs support and held-out instances from different hidden laws, and a canonical-physics baseline separate law-specific transfer from priors, appearance matching, and textbook physics. We evaluate frontier multimodal agents as complete model--harness systems and stratify accuracy by every level of the hierarchy. \TODO{Add the one-sentence quantitative headline once baseline runs finish.}
\end{abstract}
